\subsection{Relational modeling}
Our encoder-decoder approach to link prediction relies on DistMult~\cite{distmult-embedding_entities_and_relations} in the decoder,  a special and simpler case of the RESCAL factorization~\cite{nickel2011three}, more effective than the original RESCAL in the context of multi-relational knowledge bases. Numerous alternative factorizations have been proposed and studied in the context of SRL, including both (bi-)linear and nonlinear ones~(e.g., \cite{bordes2013translating,socher2013reasoning,KaiWei14,nickel2015holographic,complex-complex_embeddings_for_simple_link_prediction}). Many of these approaches can be regarded as modifications or special cases of classic tensor decomposition methods such as CP or Tucker; for a comprehensive overview of tensor decomposition literature we refer the reader to \citet{kolda2009tensor}.

Incorporation of paths between entities in knowledge
bases has recently received considerable attention. We can roughly classify previous work into (1) methods creating auxiliary triples, which are then added to the learning objective of a factorization model~\cite{guu2015traversing,garcia2015composing}; (2) approaches using paths (or walks) as features when predicting edges ~\cite{lin2015modeling}; or (3) doing both at the same time~\cite{neelakantan2015compositional,toutanova2016compositional}. The first direction is largely orthogonal to ours, as we would also expect improvements from adding similar terms to our loss (in other words, extending our decoder). The second research line is more comparable; R-GCNs provide a computationally cheaper alternative to these path-based models. %thus, letting us use larger datasets than typically considered with these methods.
Direct comparison is somewhat complicated as path-based methods used different datasets (e.g., sub-sampled sets of walks from a knowledge base).

\subsection{Neural networks on graphs}
Our R-GCN encoder model is closely related to a number of works in the area of neural networks on graphs. It is primarily motivated as an adaption of previous work on GCNs \cite{bruna2014spectral,duvenaud2015convolutional,defferrard2016convolutional,kipf2016semi} for large-scale and highly multi-relational data, characteristic of realistic knowledge bases.

Early work in this area includes the graph neural network by \citet{scarselli2009graph}. A number of extensions to the original graph neural network have been proposed, most notably \cite{li2015gated} and \cite{pham2016column}, both of which utilize gating mechanisms to facilitate optimization.

R-GCNs can further be seen as a sub-class of message passing neural networks \cite{gilmer2017neural}, which encompass a number of previous neural models for graphs, including GCNs, under a differentiable message passing interpretation.
